% Construction (v3 §1)
\section{The construction}\label{sec:construction}

The parameters $X_0$, $\kappa$, $H$, $M_1$, $L_0$, $y$, $x_k$, $\gamma_k$, $\alpha_k$, $\kappa_0$ are those of Section~\ref{sec:params}.

\subsection{The base, the rigid set and the semigroup}
\begin{itemize}
\item \emph{Base.} Below $\Xm:=X_0-H(X_0)/2$ the system is exactly $\N$. Its primes are the rational primes $p<X_0$, placed at
their own values with no jitter. On $[\Xm,X_0)$ it is $\N$ together with the at most $H(X_0)/2+1$ jittered primes that are decided
at slots $s_j<X_0+H(X_0)/2$ and land there.
\item \emph{Jittered primes.} Every prime placed by the construction (Section~\ref{sec:greedy}) has value $\ge X_0-H(X_0)/2$ and is
called \emph{jittered}.
\item \emph{Elements.} $\NN$ is the free commutative semigroup on all primes, rational and jittered. An element is an exponent
vector, and its value is $n=\prod p^{e_p}$. $N(x)$ counts elements of value $\le x$, with multiplicity. An element is
\emph{composite} if $\sum e_p\ge2$. Its \emph{latest prime} is the prime factor decided last (Section~\ref{sec:greedy}).
\item \emph{The rigid set} $\mathcal D$ is the set of $X_0$-smooth integers, i.e.\ integers all of whose prime factors are $<X_0$.
Write $\Psi(x,X_0):=\#\mathcal D\cap[1,x]$ and $\Psi(I,X_0):=\#\mathcal D\cap I$. Every $d\in\mathcal D$ with $d\ge X_0$ is composite.
\end{itemize}

\subsection{The target}\label{sec:target}
\emph{Oscillations.} For $k\ge1$ put $x_k^+:=x_k+2H(x_k)$ and
\begin{equation}\label{eq:Rk}
  R_k(x):=\alpha_kx^{\theta}\cos(\gamma_k\log x)\,\phi\bigl(\log(x/x_k^+)\bigr),\qquad R:=\sum_kR_k,
\end{equation}
where:
\begin{itemize}
\item $\phi\in C^\infty$, $0\le\phi\le1$, $\phi=0$ on $(-\infty,0]$, $\phi=1$ on $[1,\infty)$ and $\phi'\ge0$. So $R_k\equiv0$ on
$[1,x_k+2H(x_k)]$. The term $2H$ makes Lemma~\ref{lem:adapted} exact: $\gamma_k$ reads primes decided at slots up to
$x_k+H(x_k)/2$, and $T$ on $[1,u]$ is read by slots $\le u-H(u)$.
\item The \emph{amplitude condition}:
\begin{equation}\label{eq:ampl}
  \alpha_k\gamma_kx_k^{\theta-1}=\kappa'_k/\log x_k,\qquad \kappa'_k:=\kappa_0/\log^2x_k,\quad \kappa_0\in(0,\tfrac14].
\end{equation}
It gives $|R'(x)|\le\kappa_+/\log x$ with $\kappa_+=0.8$ (see \eqref{eq:G1}).
\item $\gamma_k\in[\Gamma_k,2\Gamma_k]$ is chosen in Lemma~\ref{lem:goodheights} from data on $[1,x_k)$ only, and
$\alpha_k\in[1,2]$.
\end{itemize}
In Section~\ref{sec:zeros}, $x_k$ means $x_k^+$. The difference is a relative $O(H(x_k)/x_k)$, which is invisible in every estimate.

\emph{The rigid density.} With $\kappa:=X_0^{-5/12+2\eps}$, $\eps:=0.01$, put
\begin{equation}\label{eq:rhodet}
  \rhodet(u):=\frac{\Psi([u,(1+\kappa)u],X_0)}{\kappa u}\qquad(u\ge X_0).
\end{equation}
This is a linear-window average with linear normalisation. Each rigid $n\ge X_0(1+\kappa)$ contributes a ramp of total mass
$\log(1+\kappa)/\kappa=1-\kappa/2+O(\kappa^2)$ on $[n/(1+\kappa),n]$.

\emph{The absorbed set} (random and adapted). Put $L_*:=10^7$, $h(u):=u^{\theta/2}$ and $y(n):=H(n)L_0(n)(\log n)^6$, where $H$
and $L_0$ are the parameters of Section~\ref{sec:params} at the scale of $n$. Then
\begin{equation}\label{eq:absorbed}
  \Abs:=\{n\in\NN:\ \log n\ge L_*,\ \text{the latest prime $r$ of $n$ is jittered, and } r<y(n)\},
\end{equation}
with any cofactor. Also put $\rho_{\Abs}(u):=\#(\Abs\cap[u,u+h(u)])/h(u)$.

\emph{The target.} Put $T(x):=\lfloor x\rfloor$ for $x<X_0$, and for $x\ge X_0$
\begin{equation}\label{eq:target}
\begin{gathered}
  T(x):=X_0+\bigl[\Psi(x,X_0)-\Psi(X_0,X_0)\bigr]+\#(\Abs\cap[X_0,x])+\int_{X_0}^xT'_{\mathrm{cont}}(u)\,du,\\
  T'_{\mathrm{cont}}:=1-\rhodet-\bar\rho_{\Abs}+R',\qquad \bar\rho_{\Abs}(u):=\min\bigl(\rho_{\Abs}(u),0.1/\log u\bigr).
\end{gathered}
\end{equation}
So $T$ has a unit jump at every $d\in\mathcal D\cap[X_0,\infty)$ and at every $n\in\Abs$.

\begin{remark}[The stopped process]\label{rem:stopped}
% v3: Remark 1.0
Every probabilistic statement in Sections~\ref{sec:zeros}--\ref{sec:queue} concerns the process stopped at the first failure time
$\tau$ of Definition~\ref{def:good}.
\begin{itemize}
\item Before $\tau$, i.e.\ on $\Good(<u)$, the event of Lemma~\ref{lem:lemmaP} for the window $[u,u+h(u)]$ holds, since its time is
$\le u-H(u)$ (Lemma~\ref{lem:adapted}). So $\rho_{\Abs}(u)\le(\log u)^{-4}<0.1/\log u$ and $\bar\rho_{\Abs}(u)=\rho_{\Abs}(u)$.
\item The cap can bind only after $\tau$, where nothing is claimed. Its role is to keep $T$, the slots and the construction defined on
every $\omega$.
\item The cap is adapted, being a deterministic function of the adapted $\rho_{\Abs}$.
\end{itemize}
\end{remark}

\begin{lemma}[Adaptedness]\label{lem:adapted}
% v3: Lemma 1.1
$T$ restricted to $[1,u]$ is determined by the decisions and jitters at slots $s_j\le u-H(u)$.
\end{lemma}
\begin{proof}
An element $n\le u+h(u)$ of $\Abs$ has latest prime $r<y(n)\le n^{o(1)}\ll u-H(u)$. So $r$, and every other prime of $n$, is decided
before slot $u-H(u)$ (by \eqref{eq:C2} below). The set $\mathcal D$ and the density $\rhodet$ are deterministic, and the parameters
are deterministic functions of the scale. The oscillation $R$ on $[1,u]$ involves only those $R_k$ with $x_k+2H(x_k)<u$, and each
such $\gamma_k$ reads primes decided at slots up to $x_k+H(x_k)/2\le u-H(u)$. Hence $\rho_{\Abs}$ on $[X_0,u]$, $R$ on $[1,u]$ and
the jumps of $T$ in $[1,u]$ are known before slot $u$.
\end{proof}

\subsection{Positivity of the target and the target error}
\begin{lemma}[Positivity margin]\label{lem:G1}
% v3: (G1), (1.3')
For every $u\ge X_0$ and every $\omega$,
\begin{equation}\label{eq:G1}
  T'_{\mathrm{cont}}(u)\ge\frac{0.99-0.1-\kappa_+}{\log u}-\frac{1}{\kappa u}=\frac{0.09}{\log u}-\frac1{\kappa u}>0,\qquad\kappa_+=0.8.
\end{equation}
In particular $T$ is non-decreasing on every $\omega$.
\end{lemma}
\begin{proof}
There are three inputs.
\begin{itemize}
\item \emph{Rigid part.} $1-\rhodet(u)\ge0.99/\log u-1/(\kappa u)$. Every window $[u,(1+\kappa)u]$ with $u\ge X_0$ lies in the range of
the short-interval prime number theorem (\eqref{eq:huxley}, Lemma~\ref{lem:blin}(1)), so it contains at least $0.99\kappa u/\log u$
rational primes among its at most $\kappa u+1$ integers. With the linear normalisation \eqref{eq:rhodet} there is no loss of
$\kappa/2$; such a loss belongs to the logarithmically normalised kernel, and a bound of the form $0.99/\log u-\kappa/2$ would be
negative once $\log u>2/\kappa\approx e^{3.97\cdot10^7}$, which is still below $x_1$.
\item \emph{Absorbed part.} $\bar\rho_{\Abs}\le0.1/\log u$, by the cap.
\item \emph{Oscillations.} $R\equiv0$ below $x_1+2H(x_1)$. For $x\ge x_k$, by \eqref{eq:ampl},
\[
  |R_k'(x)|\le\alpha_kx^{\theta-1}(\gamma_k+\theta+\|\phi'\|_\infty)\le2\alpha_k\gamma_kx^{\theta-1}
  =\frac{2\kappa_0(x/x_k)^{\theta-1}}{\log^3x_k}.
\]
For the last switched-on $k$: if $x\le x_k^2$ then $\log x\le2\log x_k$; otherwise $(x/x_k)^{\theta-1}\le x^{-(1-\theta)/2}$. Earlier $k$
contribute less, by the spacing $x_{k+1}\ge\exp(x_k)$. So
$|R'(x)|\le16\kappa_0/\log^3x+4\kappa_0x^{-(1-\theta)/2}\le\kappa_+/\log x$, since $\log x\ge X_0^2$.
\end{itemize}
Finally $1/(\kappa u)\le X_0^{-7/12}\ll0.09/\log u$ for $u\ge X_0$.
\end{proof}

\begin{lemma}[Target error]\label{lem:T1}
% v3: (T1), (1.4)
Write $E(x):=T(x)-x-R(x)$. On the event of Lemma~\ref{lem:lemmaP},
\begin{equation}\label{eq:T1}
  |E(x)|\le C_{X_0}x^{\theta}+2x^{\theta/2}\qquad\text{for all }x\ge X_0,\quad C_{X_0}<\infty.
\end{equation}
Moreover $|R(x)|\le2x^{\theta}\log^*x$.
\end{lemma}
\begin{proof}
\emph{Rigid part.} For $\log x\ge X_0^{1-\theta}\log X_0$, Rankin's method gives
$|\Psi(x,X_0)-\int\rhodet|\le\Psi(2x,X_0)\le(2x)^{\theta}\exp\bigl(X_0^{1-\theta}/((1-\theta)\log X_0)\bigr)$. Below that, the rigid mass
in flight gives at most $\kappa x+1$. Together this is $C_{X_0}x^{\theta}$. This includes the drift caused by the ramp mass: with the
linear normalisation, $\int_{X_0}^x\rhodet=(1-\kappa/2+O(\kappa^2))\Psi(x,X_0)$ up to the boundary terms from the rigid
$n\in[X_0,X_0(1+\kappa))$ (at most $\kappa X_0+1$, Lemma~\ref{lem:blin}(4)), so $E$ drifts by $+(\kappa/2)\Psi(x,X_0)$. At the
scales $\log x\ge X_0^2$ of the zeros this drift is $\le\kappa x^{0.011}$, since by Rankin
$\Psi(x,X_0)\le x^{0.01}\prod_{p<X_0}(1-p^{-0.01})^{-1}\le x^{0.011}$ there.

\emph{Absorbed part.} $|\#\Abs(\le x)-\int^x\rho_{\Abs}|\le\#(\Abs\cap[x-h,x+h])+O(1)\le2h(x)(\log x)^{-4}+O(1)$, by
Lemma~\ref{lem:lemmaP}.

\emph{Oscillations.} $|R(x)|\le2x^{\theta}\cdot\#\{k:x_k\le x\}\le2x^{\theta}\log^*x$.
\end{proof}

\subsection{Slots and the greedy rule}\label{sec:greedy}
\emph{Slots.} Fix $\vartheta\in(0,1)$ such that almost surely no jump of $T$ starts on a level $j+\vartheta$. This is possible: the
deterministic jump levels form a countable set, and for $n\in\Abs$ the level $T(n-)$ is a continuous strictly increasing function
of the landing position of $n$, which has a continuous law. For integers $j\ge X_0$, \emph{slot $j$} is
$J_j:=T^{-1}([j+\vartheta,j+1+\vartheta))$, with $s_j:=\inf J_j$. Each unit jump of $T$ (at $d\in\mathcal D$ or $n\in\Abs$) crosses
exactly one level and lies in its own \emph{jump slot}. Every other slot is a \emph{continuous slot}. Consecutive continuous-slot
starts are at least $1/(1+1/(\kappa u)+\kappa_+/\log u)$ apart at height $u$, since
$T'_{\mathrm{cont}}\le1+1/(\kappa u)+\kappa_+/\log u$. A single continuous slot can be shorter than $1$ where rigid integers are dense.

\emph{The hard greedy rule with box-centred jitter.} Process $j=X_0,X_0+1,\dots$ in order, keeping $Q_{j-1}\ge0$.
\begin{itemize}
\item $A_j$ is the number of composite elements whose value lies in $J_j$.
\item $P_j:=\max(0,1-Q_{j-1}-A_j)\in\{0,1\}$ and $Q_j:=Q_{j-1}+A_j+P_j-1\ge0$. No prime is placed while the queue is positive, and a
jump slot ($A_j\ge1$) never receives a prime.
\item If $P_j=1$, place one new prime at $s_j+H(s_j)(U_j-\tfrac12)$, where the $U_j$ are i.i.d.\ uniform on $[0,1]$ and independent
of the past. This holds for every slot, including those near $X_0$, so the landing density matches the slot count up to
$O(1/H)$; Section~\ref{sec:noise} uses this.
\item The \emph{jitter width} is
\begin{equation}\label{eq:jitter}
  H(x):=\max\bigl(H_{\mathrm{std}}(x),M_1(x)^2\bigr),\qquad H_{\mathrm{std}}(x):=\exp\bigl(8(\log\log x)^2\bigr),
\end{equation}
as in Section~\ref{sec:params}. \emph{Box-centred} means that the landing law of a prime decided at $s_j$ is uniform on
$[s_j-H(s_j)/2,s_j+H(s_j)/2]$.
\item \emph{Decision order.} $t(p)$ is the slot at which $p$ is decided, with $t(p):=p$ for the rational primes $p<X_0$. For an
element $w$, $\tau_w:=\max_{p\mid w}t(p)$ is the decision time of its latest prime.
\item $\mathcal F_t$ is the $\sigma$-field of all decisions and offsets at slots $\le t$.
\end{itemize}

\begin{lemma}[Causality]\label{lem:C2}
% v3: (C2)
$A_j$ is $\mathcal F_{s_j-}$-measurable:
\begin{equation}\label{eq:C2}
  A_j\in\mathcal F_{s_j-}.
\end{equation}
\end{lemma}
\begin{proof}
Let $n\in J_j$ be composite. It has a prime factor $p\le n/2<s_{j+1}/2$, which was decided at a slot $j'$ with
$s_{j'}\le p+H/2\le s_{j+1}/2+H/2<s_j$, once $s_j>2H+4$. Conversely, a prime decided at a slot $\ge j$ lands at
$\ge s_j-H/2>s_{j+1}/2$, so it is not a factor of any composite in $J_j$.
\end{proof}

\subsection{Bookkeeping}
\begin{lemma}[Count identity]\label{lem:count}
% v3: Lemma 1.2
For $j\ge X_0$ put
\[
  V_j:=\#(\text{composites of value }<s_{j+1})+\#(\text{primes decided at slots }\le j,\text{ including the rational primes }<X_0)+1.
\]
Then $V_j=j+Q_j$.
\end{lemma}
\begin{proof}
By induction: $V$ increases by $A_j+P_j$ at step $j$, and $Q_j=Q_{j-1}+A_j+P_j-1$.
\end{proof}

\begin{lemma}[Lindley]\label{lem:lindley}
% v3: Lemma 1.3
$Q_j=\max_{i\le j}\bigl[\sum_{l=i+1}^jA_l-(j-i)\bigr]$. In words, the queue is the largest excess of composite arrivals over slots,
taken over slot windows ending at $j$.
\end{lemma}
\begin{proof}
This is the standard solution of the recursion $Q_j=\max(0,Q_{j-1}+A_j-1)$, which is equivalent to the rule above.
\end{proof}

\begin{corollary}[Two-sided lag]\label{cor:lag}
% v3: Corollary 1.4, (1.5)
With $D:=N-T$, for $x\ge X_0+2H(x)$,
\begin{equation}\label{eq:lag}
  -H(x)/2-2\le D(x)\le\max_{s_j\le x}Q_j+H(x)/2+2.
\end{equation}
\end{corollary}
\begin{proof}
Landed primes $\le x$ differ from primes decided at slots $s_j\le x$ only by primes within $H/2$ of $x$, on either side. Combine this
with Lemma~\ref{lem:count}.
\end{proof}

\begin{lemma}[Excess identity]\label{lem:excess}
% v3: Lemma 1.5
Let $C_{\mathrm{jit}}(W)$ be the number of composites in $W$ whose latest prime is jittered and which are not in $\Abs$. For every slot
window $W$,
\[
  C(W)-\#\mathrm{slots}(W)=C_{\mathrm{jit}}(W)-\int_WT'_{\mathrm{cont}}\pm2.
\]
\end{lemma}
\begin{proof}
The rigid set and $\Abs$ drop out identically, because each of their elements sits on its own jump slot; the continuous slots in $W$
number $\int_WT'_{\mathrm{cont}}\pm2$. No smooth-number or Tijdeman-type input enters.
\end{proof}

\emph{Reduction.} If $C(W)-\#\mathrm{slots}(W)\le B$ for every slot window $W\subset[X_0,x]$, then $D\le B+H/2+2$ on $[X_0+2H,x]$.
Everything therefore reduces to the queue bound of Section~\ref{sec:queue} (Theorem~\ref{thm:queue}), whose drift input is
Section~\ref{sec:noise}.
